%
% Cosmic-ray air showers and their physics (chapter 2)
%
\chapter{Cosmic-Ray Air Showers}\label{chap:airshowers}
Intro, Heitler-Matthews Model.

\section{Hadronic Cascade}\label{sec:hadronic}
\subsection{Muonic Component}\label{sec:muonic}

\section{Electromagnetic Cascade}\label{sec:electromagnetic}
\subsection{Radio Emission from Air Showers}\label{sec:radio}

\section{Air Shower Universality}\label{sec:universality}
\subsection{Longitudinal Profile}\label{sec:development}

An alternative version of the Gaisser-Hillas function, is given by:

%described in terms of shower width $L$, and shower asymmetry $R$ is given by:

\begin{equation}
    N^{\prime} = \left( 1 + \frac{R X^{\prime}}{L} \right)^{R^{-2}} \exp \left( -\frac{X^{\prime}}{LR} \right). \label{eq:gh_parameterized}
\end{equation}

\noindent
$N^{\prime}$ is the normalized number of charged particles in the air shower ($N^{\prime} = N / N_{\rm max}$) and $X\prime$ is the shower maximum shifted slant depth ($X\prime = X - X_{\rm max}$).
Within this parameterization, $L$ represents a characteristic width of the shower longitudinal profile while $R$ is related to the asymmetry of the profile, defined explicitly with the equations

\begin{equation}
    L = \sqrt{|X_{0}\prime \lambda|}, \label{eq:L}
\end{equation}

\begin{equation}
    R = \sqrt{\frac{\lambda}{|X_{0}\prime|}}. \label{eq:R}
\end{equation}

\subsection{Shower-to-Shower Fluctuations}\label{sec:fluctuations}
Helium outliers from SKA proceeding and my work. Do not know where to organize it yet...

\section{Particle Attenuation in the Atmosphere}\label{sec:attenuation}
Used as motivation for including sec(zen) and r in NN

\section{Uncertainties from Hadronic Interaction Models}\label{sec:models}
